\section{Introduction}
\label{sec:intro}

Budgeted test-time inference is often summarized by one performance number: final accuracy under a
shared nominal budget of samples, branches, or reasoning actions. That number becomes ambiguous when
the nominal budget is treated as the resource measure itself and when discovery, selection, and
accounting are collapsed into a single outcome:
\begin{itemize}
\item Did any correct answer enter the candidate pool?
\item Did the commit rule select a correct answer when one was present?
\item Did compared methods share an identical pool, or did one simply receive more diverse
candidates?
\item Does a failure-derived selection rule that helps under one commercial API transfer to another?
\item Does a matched nominal budget $B$ imply comparable calls, tokens, latency, or cost?
\end{itemize}

These questions matter especially for closed-API, multi-generator evaluation: providers differ,
telemetry fields may be incomplete, and post-generation rules can appear effective until compared
to plurality on the \emph{same} extracted answers. Prior work already separates coverage from
selection and studies consensus aggregators
\citep{passatk,ll_monkeys,oracle_gap,selfconsistency}, routers, and resource reporting
\citep{snell_ttc}. The measurement gap is their
\emph{joint} operationalization under one closed-API, multi-provider, budgeted protocol with paired
identical-pool controls, transfer analysis, provenance labels, and instrumentation audit.

\begin{figure}[t]
\centering
\resizebox{\linewidth}{!}{\begin{tikzpicture}[
  font=\small,
  node distance=0.9cm and 0.55cm,
  box/.style={draw, rounded corners=1.5pt, align=center, inner sep=5pt, minimum height=1.05cm},
  wide/.style={box, text width=0.19\linewidth},
  smallbox/.style={box, text width=0.17\linewidth},
  arrow/.style={-{Latex[length=2mm]}, line width=0.45pt}
]
\node[wide] (instance) {Benchmark\\instance};
\node[wide, right=of instance] (generation) {Class A generation\\nominal budget $B$};
\node[wide, right=of generation] (pool) {Extracted candidate\\pool};
\node[wide, right=of pool] (selection) {Class B selection\\same pool, $0$ extra calls};

\draw[arrow] (instance) -- (generation);
\draw[arrow] (generation) -- (pool);
\draw[arrow] (pool) -- (selection);

\node[smallbox, below=of generation] (resources) {Realized resource\\vector $\mathbf{r}$:\\calls, tokens,\\latency, cost};
\node[smallbox, below=of pool] (discovery) {Discovery\\correct answer present?};
\node[smallbox, below=of selection] (commit) {Commit outcome\\selected answer correct?};

\draw[arrow] (generation) -- (resources);
\draw[arrow] (pool) -- (discovery);
\draw[arrow] (selection) -- (commit);

\node[box, text width=0.31\linewidth, below=1.0cm of discovery] (cells) {Provider $\times$ dataset cells:\\completed, transfer-labeled, or protocol-blocked};
\draw[arrow] (resources.south) |- (cells.west);
\draw[arrow] (commit.south) |- (cells.east);
\end{tikzpicture}
}
\caption{Protocol schematic. Nominal budget $B$ defines the logical action contract; realized
resources are measured separately. Selection claims are attributed after candidate availability is
fixed, and provider$\times$dataset cells may be completed, transfer-labeled, or protocol-blocked.}
\label{fig:protocol}
\end{figure}

\paragraph{Contribution}
This article contributes a controlled performance-evaluation protocol for budgeted closed-API
inference (Figure~\ref{fig:protocol}). The protocol combines identical-pool attribution,
provider-transfer analysis, protocol-blocked outcomes, and componentwise realized-resource
accounting under incomplete closed-API telemetry. The empirical sections use that protocol to show
four consequences: identical-pool plurality is a required selector control, failure-mined rules need
transfer labels, nominal budgets are not realized resource vectors, and provenance/blocked cells
affect interpretation even when accuracy values are unchanged.

\paragraph{Roadmap}
Section~\ref{sec:related} positions the work relative to performance measurement, test-time
compute, routing, consensus, and verifier literature. Sections~\ref{sec:problem}--\ref{sec:setup}
define the evaluation objects and closed-API workload. Sections~\ref{sec:results}--\ref{sec:compute}
report the matrix, identical-pool selection tests, provider transfer, and resource accounting.
Section~\ref{sec:ablations} reports offline counterfactuals, and
Sections~\ref{sec:limitations}--\ref{sec:conclusion} summarize scope and implications for
performance-evaluation practice.

\paragraph{Measurement objects}
The named procedures are measurement objects for exercising the protocol
(Section~\ref{sec:method}). Class~A single-generator policies under $B{=}6$ include a
reserved-attempt full-expansion diagnostic (Frontier) and closed-API adapters inspired by L1, S1,
and TALE \citep{l1,s1,tale}. Class~B post-generation selectors over the shared pool include
Pooled-4 and FTA. Offline oracles are upper bounds only. A single-cell Azure$\times$GSM8K
call/sample-budget control compares Frontier to same-model self-consistency with $N{=}6$
(SC-$N{=}6$; Section~\ref{sec:sc-entitlement-control}).
